../cictikz/src/cictikz/data/tex/cictikz_preamble.tex

% The matrix decode: a row switch on every tap, but only one column switch
% per column, so a string of 2^N taps needs 2^(N/2) + 2^(N/2) switches
% instead of the tree's 2^(N+1) - 2. The drawing is deliberately abstract —
% each switch is the knife symbol the original uses, not a device.

\begin{tikzpicture}[thick]

  \def\dx{1.6}       % column pitch
  \def\dy{1.3}       % row pitch
  \def\rr{0.20}      % switch circle radius

  % One switch: circle on the left, arm running down to the rail it drives.
  \newcommand{\rowSw}[2]{%
    \draw (#1,#2) circle (\rr);
    \draw ($(#1,#2)+(\rr,-\rr)$) -- ++(0.55,-0.45);
  }

  \foreach \c in {0,1,2,3} {
    % The column rail, from the top row down to its column switch.
    \draw ({\c*\dx+0.75}, {3*\dy+0.15}) -- ({\c*\dx+0.75}, {-0.55});
    \foreach \r in {0,1,2,3} {
      \rowSw{\c*\dx}{\r*\dy}
    }
    % The column switch, onto the shared output rail.
    \draw ({\c*\dx+0.4}, -0.9) circle (\rr);
    \draw ({\c*\dx+0.4+\rr}, {-0.9-\rr}) -- ++(0.55,-0.45);
  }

  % The output rail every column switch lands on.
  \draw (-0.35,-1.55) -- ({3*\dx+1.1}, -1.55);

  \node[anchor=east] at (-0.9, {1.5*\dy}) {Row};
  \node[anchor=north] at ({1.5*\dx+0.4}, -1.85) {Column};

\end{tikzpicture}
\end{document}
